\documentclass[10pt,a4paper]{amsart}
\usepackage{amsmath,amssymb,amsthm,graphicx}
% Retained minimal preamble for the Stacks Project chapter sites-cohomology.tex (Cohomology on Sites).
% Verbatim excerpt of the chapter preamble: Xy-pic setup (native preamble lines 14-18) and the native
% theorem declarations with the mathematical macros (native preamble lines 154-217). Copied mechanically;
% nothing added, removed or redefined. The wrapper class and the amsmath/amssymb/amsthm/graphicx base
% packages are supplied by the assembler, so the class/externaldocument part of the native preamble is omitted.
% For commutative diagrams we use Xy-pic
\usepackage[all]{xy}

% We use 2cell for 2-commutative diagrams.
\xyoption{2cell}
% Theorem environments.
%
\theoremstyle{plain}
\newtheorem{theorem}[subsection]{Theorem}
\newtheorem{proposition}[subsection]{Proposition}
\newtheorem{lemma}[subsection]{Lemma}

\theoremstyle{definition}
\newtheorem{definition}[subsection]{Definition}
\newtheorem{example}[subsection]{Example}
\newtheorem{exercise}[subsection]{Exercise}
\newtheorem{situation}[subsection]{Situation}

\theoremstyle{remark}
\newtheorem{remark}[subsection]{Remark}
\newtheorem{remarks}[subsection]{Remarks}

\numberwithin{equation}{subsection}

% Macros
%
\def\lim{\mathop{\mathrm{lim}}\nolimits}
\def\colim{\mathop{\mathrm{colim}}\nolimits}
\def\Spec{\mathop{\mathrm{Spec}}}
\def\Hom{\mathop{\mathrm{Hom}}\nolimits}
\def\Ext{\mathop{\mathrm{Ext}}\nolimits}
\def\SheafHom{\mathop{\mathcal{H}\!\mathit{om}}\nolimits}
\def\SheafExt{\mathop{\mathcal{E}\!\mathit{xt}}\nolimits}
\def\Sch{\mathit{Sch}}
\def\Mor{\mathop{\mathrm{Mor}}\nolimits}
\def\Ob{\mathop{\mathrm{Ob}}\nolimits}
\def\Sh{\mathop{\mathit{Sh}}\nolimits}
\def\NL{\mathop{N\!L}\nolimits}
\def\CH{\mathop{\mathrm{CH}}\nolimits}
\def\proetale{{pro\text{-}\acute{e}tale}}
\def\etale{{\acute{e}tale}}
\def\QCoh{\mathit{QCoh}}
\def\Ker{\mathop{\mathrm{Ker}}}
\def\Im{\mathop{\mathrm{Im}}}
\def\Coker{\mathop{\mathrm{Coker}}}
\def\Coim{\mathop{\mathrm{Coim}}}

% Boxtimes
%
\DeclareMathSymbol{\boxtimes}{\mathbin}{AMSa}{"02}

%
% Macros for moduli stacks/spaces
%
\def\QCohstack{\mathcal{QC}\!\mathit{oh}}
\def\Cohstack{\mathcal{C}\!\mathit{oh}}
\def\Spacesstack{\mathcal{S}\!\mathit{paces}}
\def\Quotfunctor{\mathrm{Quot}}
\def\Hilbfunctor{\mathrm{Hilb}}
\def\Curvesstack{\mathcal{C}\!\mathit{urves}}
\def\Polarizedstack{\mathcal{P}\!\mathit{olarized}}
\def\Complexesstack{\mathcal{C}\!\mathit{omplexes}}
% \Pic is the operator that assigns to X its picard group, usage \Pic(X)
% \Picardstack_{X/B} denotes the Picard stack of X over B
% \Picardfunctor_{X/B} denotes the Picard functor of X over B
\def\Pic{\mathop{\mathrm{Pic}}\nolimits}
\def\Picardstack{\mathcal{P}\!\mathit{ic}}
\def\Picardfunctor{\mathrm{Pic}}
\def\Deformationcategory{\mathcal{D}\!\mathit{ef}}
\usepackage{cleveref}
\crefname{originallemma}{original lemma}{original lemmas}
\Crefname{originallemma}{Original lemma}{Original lemmas}
\makeatletter
\newlabel{derived-remark-map-into-derived-limit-truncations}{{0H72}{source}{}{sourcepointer-1}{}}
\newlabel{derived-remark-map-into-derived-limit-truncations@cref}{{[originallemma][0][]0H72}{[1][0][]source}{}{}{}}
\makeatother
\crefname{originallemma}{original lemma}{original lemmas}
\Crefname{originallemma}{Original lemma}{Original lemmas}
\makeatletter
\newlabel{lemma-cohomology-derived-limit-injective}{{0D6L}{source}{}{sourcepointer-2}{}}
\newlabel{lemma-cohomology-derived-limit-injective@cref}{{[originallemma][0][]0D6L}{[1][0][]source}{}{}{}}
\makeatother
\crefname{originallemma}{original lemma}{original lemmas}
\Crefname{originallemma}{Original lemma}{Original lemmas}
\makeatletter
\newlabel{derived-remark-truncation-distinguished-triangle}{{08J5}{source}{}{sourcepointer-3}{}}
\newlabel{derived-remark-truncation-distinguished-triangle@cref}{{[originallemma][0][]08J5}{[1][0][]source}{}{}{}}
\makeatother
\title{An object with vanishing cohomology of its cohomology sheaves over a cofinal system is the derived limit of its truncations}
\author{The Stacks Project Authors}
\date{stacks-project:sites-cohomology.tex:local-snapshot-20261011; GNU FDL 1.2 or later}
\begin{document}
\maketitle
\noindent\textit{Editorial note.} All delivered mathematical text is the author's own, copied from the pinned Stacks Project chapter file sites-cohomology.tex; no statement or proof step was rewritten, reordered or removed.\par\medskip
\noindent\textit{Editorial note.} The proof is the author's complete argument; the reduction to the cohomology sheaves and the use of the vanishing criterion are both inside the counted span.\par\medskip
\noindent\textit{Editorial note.} Every labelled reference occurring inside the counted spans is carried in assembly.external\_references as a non-load-bearing navigation pointer (same-chapter ones at their definition site, cross-chapter ones at the invocation inside the counted span), so no load-bearing mathematics is replaced by a stub. No bibliography citation occurs inside the counted spans.\par\medskip
\noindent\textit{Editorial note.} The checked chapter list for same-name collisions is the set of stacks chapters staged in this repository: chow, cohomology, divisors, etale-cohomology, flat, more-algebra, more-morphisms, morphisms, spaces-more-morphisms, varieties together with this card's chapter sites-cohomology.\par\medskip
\noindent\textit{Editorial note.} The preamble is the chapter's own retained minimal preamble (sources/stacks-sites-cohomology/retained-preamble.tex), a verbatim excerpt of the native preamble: the Xy-pic setup and the native theorem declarations with the mathematical macros.\par\medskip
\noindent\textit{Allowed standard prerequisites (editorial).}\begin{itemize}
\item Truncation functors tau\_\{>= -n\} and the derived limit of a system of objects. Same chapter together with the chapter Derived Categories, the truncations of an object of D(O) and the construction of R lim of an inverse system of complexes. Supplies the object R lim tau\_\{>= -n\}E and the canonical map from E to it whose being an isomorphism is asserted.
\item Cohomology of a complex of O-modules at an object of the site, computed over a cofinal system of coverings. Same chapter, the vanishing criterion for cohomology of sheaves over a cofinal system of coverings, which is the hypothesis (2) of the statement. Supplies the input that makes the successive quotients of the limit vanish.
\end{itemize}\medskip
\ifdefined\hypertarget\hypertarget{sourcepointer-1}{}\fi
\noindent\textit{Source reference.} Original lemma 0H72: sources/stacks-sites-cohomology/pinned-original-source.tex, source lines 5337--5337. Navigation pointer invoked at proof line 5337; the counted span contains the pointer only and the cited result is not reproduced inside it, so the pointer is registered as a non-load-bearing navigation stub in the established form for Stacks material. No bibliography citation is involved.\par\medskip
\ifdefined\hypertarget\hypertarget{sourcepointer-2}{}\fi
\noindent\textit{Source reference.} Original lemma 0D6L: sources/stacks-sites-cohomology/pinned-original-source.tex, source lines 5279--5279. Navigation pointer invoked at proof line 5371; the counted span contains the pointer only and the cited result is not reproduced inside it, so the pointer is registered as a non-load-bearing navigation stub in the established form for Stacks material. No bibliography citation is involved.\par\medskip
\ifdefined\hypertarget\hypertarget{sourcepointer-3}{}\fi
\noindent\textit{Source reference.} Original lemma 08J5: sources/stacks-sites-cohomology/pinned-original-source.tex, source lines 5374--5374. Navigation pointer invoked at proof line 5374; the counted span contains the pointer only and the cited result is not reproduced inside it, so the pointer is registered as a non-load-bearing navigation stub in the established form for Stacks material. No bibliography citation is involved.\par\medskip
\begin{lemma}
\label{lemma-is-limit-per-object}
Let $(\mathcal{C}, \mathcal{O})$ be a ringed site. Let $E \in D(\mathcal{O})$.
Let $\mathcal{B} \subset \Ob(\mathcal{C})$ be a subset. Assume
\begin{enumerate}
\item every object of $\mathcal{C}$ has a covering whose members
are elements of $\mathcal{B}$, and
\item for every $V \in \mathcal{B}$ there exist a function
$p(V, -) : \mathbf{Z} \to \mathbf{Z}$ and a cofinal system $\text{Cov}_V$
of coverings of $V$ such that
$$
H^p(V_i, H^{m - p}(E)) = 0
$$
for all $\{V_i \to V\} \in \text{Cov}_V$ and all integers $p, m$
satisfying $p > p(V, m)$.
\end{enumerate}
Then the map $E \to R\lim \tau_{\geq -n} E$ of
Derived Categories, Remark
\ref{derived-remark-map-into-derived-limit-truncations}
is an isomorphism in $D(\mathcal{O})$.
\end{lemma}

\begin{proof}
Set $K_n = \tau_{\geq -n}E$ and $K = R\lim K_n$.
The canonical map $E \to K$
comes from the canonical maps $E \to K_n = \tau_{\geq -n}E$.
We have to show that $E \to K$ induces an isomorphism
$H^m(E) \to H^m(K)$ of cohomology sheaves. In the rest of the
proof we fix $m$. If $n \geq -m$, then
the map $E \to \tau_{\geq -n}E = K_n$ induces an isomorphism
$H^m(E) \to H^m(K_n)$.
To finish the proof it suffices to show that for every $V \in \mathcal{B}$
there exists an integer $n(V) \geq -m$ such that the map
$H^m(K)(V) \to H^m(K_{n(V)})(V)$ is injective. Namely, then
the composition
$$
H^m(E)(V) \to H^m(K)(V) \to H^m(K_{n(V)})(V)
$$
is a bijection and the second arrow is injective, hence the
first arrow is bijective. By property (1) this will imply
$H^m(E) \to H^m(K)$ is an isomorphism. Set
$$
n(V) = 1 + \max\{-m, p(V, m - 1) - m, -1 + p(V, m) - m, -2 + p(V, m + 1) - m\}.
$$
so that in any case $n(V) \geq -m$. Claim: the maps
$$
H^{m - 1}(V_i, K_{n + 1}) \to H^{m - 1}(V_i, K_n)
\quad\text{and}\quad
H^m(V_i, K_{n + 1}) \to H^m(V_i, K_n)
$$
are isomorphisms for $n \geq n(V)$ and $\{V_i \to V\} \in \text{Cov}_V$.
The claim implies conditions
(1) and (2) of Lemma \ref{lemma-cohomology-derived-limit-injective}
are satisfied and hence implies the desired injectivity.
Recall (Derived Categories, Remark
\ref{derived-remark-truncation-distinguished-triangle})
that we have distinguished triangles
$$
H^{-n - 1}(E)[n + 1] \to
K_{n + 1} \to K_n \to H^{-n - 1}(E)[n + 2]
$$
Looking at the associated long exact cohomology sequence the claim follows if
$$
H^{m + n}(V_i, H^{-n - 1}(E)),\quad
H^{m + n + 1}(V_i, H^{-n - 1}(E)),\quad
H^{m + n + 2}(V_i, H^{-n - 1}(E))
$$
are zero for $n \geq n(V)$ and $\{V_i \to V\} \in \text{Cov}_V$.
This follows from our choice of $n(V)$
and the assumption in the lemma.
\end{proof}

\end{document}
